\section{Methodology}
\label{sec:method}

\begin{figure}[t]
    \centering
    \resizebox{0.98\linewidth}{!}{\begin{tikzpicture}[
    box/.style={draw, rounded corners=2pt, align=center, font=\small, inner sep=5pt, text width=#1},
    box/.default=2.7cm,
    arr/.style={-{Latex[length=2mm]}, thick},
    node distance=0.55cm]
  \node[box] (units) {\textbf{44 units}\\4 models $\times$\\11 interventions\\(4 are placebos)};
  \node[box, right=of units] (out) {\textbf{Model outputs}\\240 shared items,\\baseline and\\intervention};
  \node[box=3.4cm, right=of out, yshift=1.3cm] (multi) {\textbf{Multiverse}\\4 metrics $\times$ 4 sample sizes;\\1000 resampled studies per cell\\$\rightarrow$ null rate (704 cells)};
  \node[box=3.4cm, right=of out, yshift=-1.3cm] (ref) {\textbf{Reference effect}\\all 240 items, Holm-corrected\\$\rightarrow$ \behav, \latent\\or \noeff};
  \node[box=2.6cm, right=of multi, yshift=-1.3cm] (q) {\textbf{Questions}\\Where do nulls\\cluster?\\What are they\\made of?};
  \draw[arr] (units) -- (out);
  \draw[arr] (out.east) -- (multi.west);
  \draw[arr] (out.east) -- (ref.west);
  \draw[arr] (multi.east) -- (q.west);
  \draw[arr] (ref.east) -- (q.west);
\end{tikzpicture}
}
    \caption{Overview of the design. Each model--intervention pair (a \emph{unit}) is run once on 240 shared items. From the recorded outputs we derive two things: a multiverse of resampled studies that gives a null rate for every metric and sample size, and a reference effect on the full pool that classifies the unit. Comparing the two tells us where nulls cluster and what kind of underlying effect they sit on.}
    \label{fig:overview}
\end{figure}

\subsection{Problem formulation}
\label{sec:formulation}

Let a \emph{unit} $u = (m, i)$ be a model $m$ paired with an intervention $i$.
For each evaluation item $j$ and metric $k$, we record a per-item score $y^{k}_{u,j}$ under the intervention and $y^{k}_{m,j}$ at baseline.
A \emph{study} is a tuple $(u, k, n)$: it draws $n$ items, runs a paired two-sided test of intervention against baseline on metric $k$, and reports a \emph{null} if $p \ge \alpha$, with $\alpha = .05$.
The \emph{null rate} of a design cell $(u, k, n)$ is the fraction of studies in that cell that report a null.
Our main question is how the null rate depends on four design features: metric $k$, sample size $n$, the granularity of intervention $i$, and the scale of model $m$.

Metric and sample size are \emph{measurement-side} features.
Within a unit, the model outputs are identical across their levels, so a difference in null rate between metrics is methodological by construction.
Granularity and scale are \emph{effect-side} features: changing them changes the intervention or the model, so a difference in null rate may reflect a truly different effect.
We fixed this asymmetry in the plan before collecting data, and it is why we test the ``blind spot'' interpretation separately (H4 below).

We test four hypotheses.
{\bf H1 (metric):} holding model, intervention, and items fixed, null rate differs by metric, with answer-level metrics giving more nulls than continuous ones.
{\bf H2 (granularity):} null rate differs by intervention granularity.
{\bf H3 (scale):} null rate differs by model size.
{\bf H4 (blind spot):} among nulls from non-placebo interventions, the majority occur where the reference analysis shows a real behavioral effect.

\subsection{Models, data, and prompt}

\para{Models.}
We use \qwen at 0.5B, 1.5B, 3B, and 7B parameters~\cite{qwen2024qwen25}, run locally on one NVIDIA RTX A6000 with batch size 8.
We choose open-weight models because three of our four intervention types need access to activations.
The three smaller models run in float32.
The 7B model runs in bfloat16 with the unembedding computed in float32, because the GPUs were shared.

\para{Data.}
We use the sycophancy sets from the model-written evaluations of \citet{perez2022discovering}: NLP survey, PhilPapers 2020, and political typology, restricted to binary-choice items.
Each item contains a user biography that implies a view, followed by a two-option question; one option matches the user's implied view.
We draw a seeded sample balanced on which letter is the sycophantic one: 32 train and 80 evaluation items per domain, disjoint.
The 96 train items are used only to build interventions.
The 240 evaluation items are identical across all models and conditions.

\para{Prompt.}
We use the default \qwen system prompt, the item as the user turn, and an assistant turn prefilled with `` (\,'' so that the next token is the answer letter.
One forward pass per item gives the full next-token distribution, from which we compute all metrics.
Mean probability mass on the two answer letters is at least 0.93 in every condition except \intv{S\_0.3} on 7B, where it is 0.90.

\subsection{Interventions}

\begin{table}[t]
    \centering
    \caption{The eleven interventions, grouped by granularity. Each granularity has one placebo ($\dagger$) that matches the form of the real intervention but is not aimed at sycophancy. Layer, doses, and head counts were fixed in advance and not tuned on evaluation items.}
    \label{tab:interventions}
    \small
    \begin{tabular}{@{}llp{8.6cm}@{}}
        \toprule
        \textbf{Granularity} & \textbf{Name} & \textbf{What it does} \\
        \midrule
        Prompt & \intv{P\_system} & System instruction not to tailor answers to the user's views \\
               & \intv{P\_user} & User-turn request to ignore the stated views \\
               & \intv{P\_placebo}$^\dagger$ & Irrelevant formatting instruction \\
        \midrule
        All layers & \intv{D\_ablate} & Project the sycophancy direction out after every block \\
                   & \intv{D\_ablate\_rand}$^\dagger$ & Same with a random direction \\
        \midrule
        Single layer & \intv{S\_0.1}, \intv{S\_0.3} & Subtract the direction at the middle layer, scaled to 0.1 or 0.3 of the mean residual norm \\
                     & \intv{S\_rand\_0.3}$^\dagger$ & Same with a random direction at dose 0.3 \\
        \midrule
        Head & \intv{H\_top1}, \intv{H\_topk} & Zero-ablate the top 1 head or the top 2\% of heads by attribution \\
             & \intv{H\_randk}$^\dagger$ & Zero-ablate the same number of random heads \\
        \bottomrule
    \end{tabular}
\end{table}

\Tabref{tab:interventions} lists the interventions.
We choose four granularities because they span the range used in practice, from editing the prompt to editing a handful of attention heads.

\para{Sycophancy direction.}
Following the contrastive approach of \citet{rimsky2023steering}, we append the sycophantic or the non-sycophantic letter to each train prompt, read the residual stream after the middle block (index $L/2$ for a model with $L$ blocks) at the letter position, and take the mean difference over the 96 train items, normalized to unit length.
\intv{D\_ablate} projects this direction out of the output of every block, in the style of \citet{arditi2024refusal}.
\intv{S\_0.1} and \intv{S\_0.3} subtract the direction at the middle block at every position, scaled to 0.1 or 0.3 of the mean residual norm at that block.

\para{Head selection.}
We score each attention head by gradient $\times$ activation attribution to the sycophancy logit margin, using 48 of the train items, and zero-ablate the top-scoring head (\intv{H\_top1}) or the top 2\% of heads (\intv{H\_topk}; 7, 7, 12, and 16 heads for the four models).
\intv{H\_randk} ablates the same number of heads chosen at random.

\subsection{Metrics}

We compute four per-item metrics from the same next-token distribution.
Let $s$ be the sycophantic letter and $\bar{s}$ the other one.
\begin{itemize}[leftmargin=*,itemsep=0pt,topsep=0pt]
    \item \msample: one draw at temperature 1 equals $s$ (binary).
    \item \mgreedy: the most likely token equals $s$ (binary).
    \item \mprob: $P(s) / (P(\texttt{A}) + P(\texttt{B}))$.
    \item \mlogit: the logit margin, $\mathrm{logit}(s) - \mathrm{logit}(\bar{s})$.
\end{itemize}
The first two are \emph{answer-level} metrics: they record what the model says.
The last two are continuous and record how strongly it prefers one answer.
All four are in common use, and nothing about the experiment changes between them except what is written down.

\subsection{Multiverse of studies}

\begin{algorithm}[t]
\caption{Null rate of one design cell}
\label{alg:multiverse}
\begin{algorithmic}[1]
\Require unit $u$, metric $k$, sample size $n$, per-item scores for 240 items, $R = 1000$
\State $c \gets 0$
\For{$r = 1, \dots, R$}
    \State Draw $n$ items with replacement from the 240-item pool
    \If{$k = \msample$}
        \State Draw fresh Bernoulli answers from the recorded $P(s)$, for baseline and intervention
    \EndIf
    \State $p \gets$ exact McNemar test if $k$ is binary, paired $t$-test otherwise \Comment{two-sided}
    \If{$p \ge .05$}
        \State $c \gets c + 1$ \Comment{this study reports a null}
    \EndIf
\EndFor
\State \Return $c / R$
\end{algorithmic}
\end{algorithm}

Algorithm~\ref{alg:multiverse} gives the procedure for one cell.
We use $n \in \{20, 50, 100, 200\}$, which covers the range from a quick check to a moderately sized evaluation.
With 44 units, 4 metrics, and 4 sample sizes, the map has 704 cells and 704,000 simulated studies.
The model outputs are real; only the sampling of items into studies is simulated.

\subsection{Reference effects}
\label{sec:reference}

To interpret a null we need to know what the intervention does when measured as well as we can.
For each unit we test the intervention against baseline on all 240 items, using a paired $t$-test for the logit margin and an exact McNemar test for the greedy answer, with Holm correction across the 44 units.
We then assign one of three classes:
\begin{itemize}[leftmargin=*,itemsep=0pt,topsep=0pt]
    \item \behav: the greedy sycophantic-answer rate changes.
    \item \latent: the logit margin changes, but the greedy rate does not.
    \item \noeff: neither changes.
\end{itemize}
We pre-specified that nulls in \behav units count as missed effects, and that nulls in \latent units are reported separately, because a logit shift that flips no answers is arguably robustness at the behavioral level.

\subsection{Statistical analysis}

Cells that share a unit are not independent, so we use resampling methods that respect the unit structure.
For metric and sample size we use within-unit permutation tests (5000 permutations) with the range of marginal null rates as the statistic.
For granularity we permute intervention labels.
For model scale we report descriptive results only, because four models from one family do not support an inferential claim.
Confidence intervals come from a cluster bootstrap over the 44 units (2000 resamples).
We report shares of variance from a type-II ANOVA on cell null rates as a descriptive summary and do not rely on its $F$-tests.
Analyses that we added after seeing the results are marked as exploratory.
